% =============================================================================
%  progress-report-03.tex — รายงานความก้าวหน้าโครงการ ULMs ครั้งที่ 3
%  *** ฉบับส่งอาจารย์ รอบ 4 – 10 ก.ย. 2569 (วันสิ้นสุด Sprint 3) ***
%
%  build:  cd ~/Projects/SoftwareEn/docs && latexmk progress-report-03.tex
%          ผลลัพธ์ออกที่ build/progress-report-03.pdf
%
%  หลักการเขียน — ยกมาจาก progress-report-02.tex ทุกข้อ อ่านหัวไฟล์นั้นก่อนแก้
%    สั้น ไม่มีสารบัญ · แผนเส้นเดียว · อ้างได้เฉพาะฉบับที่ส่งจริง (ครั้งที่ 1 และ 2)
%    ห้ามรายงานเป็นจำนวนบรรทัด · ไม่มี \textbf กลางย่อหน้า
%    ศัพท์ของการพัฒนาใช้อังกฤษ และเว้นวรรคหน้าหลังทุกครั้งที่ชนอักษรไทย
%    ห้ามรัน regex ยุบช่องไฟทับทั้งไฟล์
%
%  ห้ามแก้ progress-report-02.tex ย้อนหลัง ฉบับนั้นส่งไปแล้ว
%
%  ตัวเลขทุกตัวอ้าง origin/main ที่ c8ed505 (10 ก.ย. 2569 09:22)
%  fetch ซ้ำก่อนสร้าง PDF แล้ว ไม่มี commit ใหม่ ไม่นับ spike-option-b/
%
%  วันรายงานตรงกับขอบ Sprint 3 พอดี ค่าแผนจึงอ่านตรงจากแถว S3 ของภาคผนวก
%  ไม่ต้องสอดแทรกเหมือนรอบก่อน
%
%  ที่มาของตัวเลข ตรวจซ้ำได้จาก repo:
%    86 procedure / 10 domain = ผลรวม @Query/@Mutation ใน backend/src/*/*.router.ts
%    30 ตาราง 9 enum 9 migration = backend/prisma/
%    234 test      = เลขที่ jest รายงานจาก `npx jest` ในเครื่องผ่าน 208 ตก 26
%                    เพราะไม่มี Postgres (ตกชุดเดิมกับรอบก่อน)
%    21 spec file; 8 จาก 21 service มี .spec.ts คู่กัน (ไม่นับ app กับ prisma
%                    ตามวิธีนับของรอบก่อนที่ได้ 18)
%    44 AC check   = `npm run test:ac` ผ่านทั้งหมด
%    28 หน้า       = frontend/src/features/**/*-page.tsx (รอบก่อน 30 ลบไป 2)
%    26 มี UI จริง = ไม่ใช้ PlaceholderPage (เหลือ handover กับ reports/export)
%    21 อ่านข้อมูลจริง / 5 ข้อมูลหลักยังเป็น mock (rooms 3, appeals 2)
%    11 procedure ที่ FE ประกาศแต่ไม่มีจริง = FE 74 (11 domain) เทียบ BE 86
%    รหัสข้อผิดพลาดที่มีคำแปลไทย 60 จาก 63 (รอบก่อน 29 จาก 59)
%    21 commit / 9 PR ในช่วง 99da590..c8ed505 (ไม่รวม merge = 11)
% =============================================================================
% =============================================================================
%  ulms-preamble.tex — preamble ที่ใช้ร่วมกันทุกเอกสารของโครงการ ULMs
%
%  ใช้โดย:  main.tex (ข้อเสนอโครงการ) · progress.tex (รายงานความก้าวหน้า)
%
%  เอกสารที่เรียกใช้ต้องทำสองอย่างหลัง % =============================================================================
%  ulms-preamble.tex — preamble ที่ใช้ร่วมกันทุกเอกสารของโครงการ ULMs
%
%  ใช้โดย:  main.tex (ข้อเสนอโครงการ) · progress.tex (รายงานความก้าวหน้า)
%
%  เอกสารที่เรียกใช้ต้องทำสองอย่างหลัง % =============================================================================
%  ulms-preamble.tex — preamble ที่ใช้ร่วมกันทุกเอกสารของโครงการ ULMs
%
%  ใช้โดย:  main.tex (ข้อเสนอโครงการ) · progress.tex (รายงานความก้าวหน้า)
%
%  เอกสารที่เรียกใช้ต้องทำสองอย่างหลัง \input{ulms-preamble}:
%    1) \hypersetup{pdftitle={...}}  — ชื่อเอกสารใน metadata ของ PDF
%    2) \begin{document}
%  (pdftitle ถูกย้ายออกจากไฟล์นี้ เพราะเป็นค่าเฉพาะของแต่ละเอกสาร)
%
%  แก้ที่นี่ที่เดียว = มีผลกับทุกเอกสาร ห้ามก๊อป preamble ไปวางซ้ำ
% =============================================================================
% =============================================================================
%  main.tex  —  เอกสารข้อเสนอโครงการ (Academic Project Proposal)
%  ระบบบริหารจัดการการยืม–คืนอุปกรณ์มหาวิทยาลัย (ULMs)
%
%  IMPORTANT: This document contains Thai text and MUST be compiled with
%  XeLaTeX (not pdflatex). Thai needs a Unicode font + ICU line-breaking.
%      Build:  latexmk main.tex        (uses .latexmkrc -> xelatex)
%      or:     xelatex main.tex        (run twice for the table of contents)
% =============================================================================

\documentclass[11pt, a4paper]{article}

% ---- Fonts & Thai support ----------------------------------------------------
% fontspec lets XeLaTeX use any system font. Laksaman is a complete serif face
% that covers BOTH Thai and Latin glyphs, so one \setmainfont handles the whole
% (Thai + English) document without font switching.
\usepackage{fontspec}

% Thai has no spaces between words, so TeX cannot find line-break points on its
% own. These two XeTeX primitives switch on ICU's dictionary-based Thai word
% segmentation, which is what makes Thai paragraphs wrap correctly.
\XeTeXlinebreaklocale "th"
\XeTeXlinebreakskip = 0pt plus 0.1pt

\setmainfont{Laksaman}

% ---- Page geometry & paragraph style ----------------------------------------
\usepackage[margin=2.5cm]{geometry}
\usepackage{parskip}          % block paragraphs (no indent, gap between) — common
                              % and more readable for Thai documents
\linespread{1.15}             % a little extra leading; Thai has tall tone marks

% ---- Colours (monochrome / grayscale palette) ------------------------------
\usepackage[table]{xcolor}    % 'table' option enables \rowcolor for table headers
% Monochrome (grayscale) palette. Names kept as-is so every reference still
% works; only the tones changed from the original green/blue accents.
\definecolor{kugreen}{HTML}{1A1A1A}   % near-black — section headings & title
\definecolor{kudark}{HTML}{333333}    % dark gray — subsection/subsubsection titles
\definecolor{headbg}{HTML}{EAEAEA}    % light gray background for table header rows
\definecolor{linkblue}{HTML}{404040}  % medium-dark gray — links / citations

% ---- Section heading styling -------------------------------------------------
% 'nobottomtitles*' stops a heading from being printed near the bottom of a
% page: if less than \bottomtitlespace of room is left, the heading (and its
% following text) moves to the next page. This is titlesec's built-in, robust
% replacement for the earlier \needspace hack — the latter injected vertical
% glue that clashed with longtable page-splitting ("Infinite glue shrinkage").
\usepackage[nobottomtitles*]{titlesec}
\titleformat{\section}
  {\color{kugreen}\Large\bfseries}{\thesection}{0.6em}{}
\titleformat{\subsection}
  {\color{kudark}\large\bfseries}{\thesubsection}{0.5em}{}
\titleformat{\subsubsection}
  {\color{kudark}\normalsize\bfseries}{\thesubsubsection}{0.5em}{}
\titlespacing*{\section}{0pt}{1.4\baselineskip}{0.5\baselineskip}

% Minimum room a heading needs at the bottom of a page; if only ~1–3 lines are
% left (less than this), the heading is bumped to the next page instead of being
% stranded above its table/text.
\renewcommand{\bottomtitlespace}{4\baselineskip}

% ---- Hierarchical indentation by heading level ------------------------------
% Each heading level shifts BOTH the heading and the body text under it further
% to the right: section = 0, subsection (x.y) = 1 step, subsubsection (x.y.z) =
% 2 steps. We do this by having the sectioning commands set \leftskip (a global
% left margin that carries through every following line until the next heading
% resets it). \par first, so the *previous* paragraph is not re-indented.
% \leftskip indents plain paragraphs; \@totalleftmargin makes list environments
% (itemize/enumerate) start from the same indent instead of the page margin.
\makeatletter
\newlength{\lvlindent}
\setlength{\lvlindent}{1.8em}          % one indent step
\let\ulmsSecIndent\section
\renewcommand{\section}{\par\global\leftskip=0pt\global\@totalleftmargin=0pt\relax\ulmsSecIndent}
\let\ulmsSubsecIndent\subsection
\renewcommand{\subsection}{\par\global\leftskip=\lvlindent\global\@totalleftmargin=\lvlindent\relax\ulmsSubsecIndent}
\let\ulmsSubsubsecIndent\subsubsection
\renewcommand{\subsubsection}{\par\global\leftskip=2\lvlindent\global\@totalleftmargin=2\lvlindent\relax\ulmsSubsubsecIndent}
\makeatother

% Indent the headings themselves to match (titlesec ignores \leftskip for the
% heading line, so its left margin is set here via the first \titlespacing arg).
\titlespacing*{\subsection}   {\lvlindent} {1.0\baselineskip}{0.4\baselineskip}
\titlespacing*{\subsubsection}{2\lvlindent}{0.8\baselineskip}{0.3\baselineskip}

% ---- Lists (tighter, custom bullets) ----------------------------------------
\usepackage{enumitem}
\setlist[itemize]{leftmargin=1.6em, itemsep=2pt, topsep=3pt}
\setlist[enumerate]{leftmargin=1.9em, itemsep=2pt, topsep=3pt}

% ---- Tables ------------------------------------------------------------------
\usepackage{array}
\usepackage{booktabs}         % \toprule / \midrule / \bottomrule
\usepackage{tabularx}         % auto-width columns that wrap long Thai text
\usepackage{longtable}        % tables that can break across pages
\usepackage{multirow}         % merged cells for the team-structure table

% Left-aligned wrapping column types:
%   L{width} — fixed-width ragged-right paragraph column
%   Y        — flexible ragged-right X column (for tabularx)
\newcolumntype{L}[1]{>{\raggedright\arraybackslash}p{#1}}
\newcolumntype{Y}{>{\raggedright\arraybackslash}X}
\renewcommand{\arraystretch}{1.3}

% Helpers for a styled header row inside tables.
%   \hrow  — starts a header row and shades it (must be FIRST in the row;
%            colortbl allows only one \rowcolor per row).
%   \thead — bold header-cell text.
\newcommand{\hrow}{\rowcolor{headbg}}
\newcommand{\thead}[1]{\textbf{#1}}

% \begin{keeptable} … \end{keeptable} — keeps a table and the bold label above
% it on the SAME page (used by the data dictionary in §5). A tabularx cannot
% split across pages, so LaTeX is free to break between the label and its
% table, stranding the label alone at the bottom of a page; wrapping the pair
% in a minipage makes them one unbreakable block.
%   * \leftskip is zeroed inside: the minipage box is already placed at the
%     current heading indent, so leaving it set would indent the content twice
%     and push the table past the right margin.
%   * \parskip is restored because minipage's \@parboxrestore zeroes it, which
%     would glue the table straight onto the label line.
% Tables inside must use \linewidth (= the minipage width), not \textwidth.
\newenvironment{keeptable}
  {\par\medskip\noindent
   \begin{minipage}{\dimexpr\textwidth-\leftskip\relax}%
   \leftskip=0pt\relax
   \setlength{\parskip}{0.5\baselineskip plus 2pt}}
  {\end{minipage}\par\medskip}

% ---- Table of contents styling ----------------------------------------------
% By default LaTeX gives dotted leaders to subsections but NOT to sections
% (they print bold with just a right-aligned page number). tocloft lets us turn
% on dots for the section level too, so every TOC line runs "….." to its page.
\usepackage{tocloft}
\renewcommand{\cftsecdotsep}{\cftdotsep}   % dotted leaders for section entries
\renewcommand{\cftsecleader}{\cftdotfill{\cftsecdotsep}}
\renewcommand{\cftsecpagefont}{\normalfont}  % keep page numbers un-bold

% ---- Figures -----------------------------------------------------------------
\usepackage{graphicx}
\graphicspath{{figures/}}

% Landscape (rotated) full-page floats — needed by the ER diagram in §5, which
% is far wider than it is tall and would be unreadable scaled to \textwidth.
\usepackage{rotating}

% ---- Flowcharts / diagrams (TikZ) -------------------------------------------
% Used for the borrow–return workflow diagram (§5). Libraries: geometric shapes
% for decision diamonds, arrows.meta for arrowheads, positioning for relative
% node placement, calc/fit/backgrounds for routing loops and grouping.
\usepackage{tikz}
\usetikzlibrary{shapes.geometric, arrows.meta, positioning, calc, fit, backgrounds}

% ---- Links & bookmarks -------------------------------------------------------
\usepackage{hyperref}
\hypersetup{
  colorlinks=true, linkcolor=black, citecolor=linkblue, urlcolor=linkblue,
  pdflang={th},
}

% Thai labels for auto-generated headings.
\renewcommand{\contentsname}{สารบัญ (Table of Contents)}
\renewcommand{\tablename}{ตาราง}
\renewcommand{\figurename}{รูปที่}

:
%    1) \hypersetup{pdftitle={...}}  — ชื่อเอกสารใน metadata ของ PDF
%    2) \begin{document}
%  (pdftitle ถูกย้ายออกจากไฟล์นี้ เพราะเป็นค่าเฉพาะของแต่ละเอกสาร)
%
%  แก้ที่นี่ที่เดียว = มีผลกับทุกเอกสาร ห้ามก๊อป preamble ไปวางซ้ำ
% =============================================================================
% =============================================================================
%  main.tex  —  เอกสารข้อเสนอโครงการ (Academic Project Proposal)
%  ระบบบริหารจัดการการยืม–คืนอุปกรณ์มหาวิทยาลัย (ULMs)
%
%  IMPORTANT: This document contains Thai text and MUST be compiled with
%  XeLaTeX (not pdflatex). Thai needs a Unicode font + ICU line-breaking.
%      Build:  latexmk main.tex        (uses .latexmkrc -> xelatex)
%      or:     xelatex main.tex        (run twice for the table of contents)
% =============================================================================

\documentclass[11pt, a4paper]{article}

% ---- Fonts & Thai support ----------------------------------------------------
% fontspec lets XeLaTeX use any system font. Laksaman is a complete serif face
% that covers BOTH Thai and Latin glyphs, so one \setmainfont handles the whole
% (Thai + English) document without font switching.
\usepackage{fontspec}

% Thai has no spaces between words, so TeX cannot find line-break points on its
% own. These two XeTeX primitives switch on ICU's dictionary-based Thai word
% segmentation, which is what makes Thai paragraphs wrap correctly.
\XeTeXlinebreaklocale "th"
\XeTeXlinebreakskip = 0pt plus 0.1pt

\setmainfont{Laksaman}

% ---- Page geometry & paragraph style ----------------------------------------
\usepackage[margin=2.5cm]{geometry}
\usepackage{parskip}          % block paragraphs (no indent, gap between) — common
                              % and more readable for Thai documents
\linespread{1.15}             % a little extra leading; Thai has tall tone marks

% ---- Colours (monochrome / grayscale palette) ------------------------------
\usepackage[table]{xcolor}    % 'table' option enables \rowcolor for table headers
% Monochrome (grayscale) palette. Names kept as-is so every reference still
% works; only the tones changed from the original green/blue accents.
\definecolor{kugreen}{HTML}{1A1A1A}   % near-black — section headings & title
\definecolor{kudark}{HTML}{333333}    % dark gray — subsection/subsubsection titles
\definecolor{headbg}{HTML}{EAEAEA}    % light gray background for table header rows
\definecolor{linkblue}{HTML}{404040}  % medium-dark gray — links / citations

% ---- Section heading styling -------------------------------------------------
% 'nobottomtitles*' stops a heading from being printed near the bottom of a
% page: if less than \bottomtitlespace of room is left, the heading (and its
% following text) moves to the next page. This is titlesec's built-in, robust
% replacement for the earlier \needspace hack — the latter injected vertical
% glue that clashed with longtable page-splitting ("Infinite glue shrinkage").
\usepackage[nobottomtitles*]{titlesec}
\titleformat{\section}
  {\color{kugreen}\Large\bfseries}{\thesection}{0.6em}{}
\titleformat{\subsection}
  {\color{kudark}\large\bfseries}{\thesubsection}{0.5em}{}
\titleformat{\subsubsection}
  {\color{kudark}\normalsize\bfseries}{\thesubsubsection}{0.5em}{}
\titlespacing*{\section}{0pt}{1.4\baselineskip}{0.5\baselineskip}

% Minimum room a heading needs at the bottom of a page; if only ~1–3 lines are
% left (less than this), the heading is bumped to the next page instead of being
% stranded above its table/text.
\renewcommand{\bottomtitlespace}{4\baselineskip}

% ---- Hierarchical indentation by heading level ------------------------------
% Each heading level shifts BOTH the heading and the body text under it further
% to the right: section = 0, subsection (x.y) = 1 step, subsubsection (x.y.z) =
% 2 steps. We do this by having the sectioning commands set \leftskip (a global
% left margin that carries through every following line until the next heading
% resets it). \par first, so the *previous* paragraph is not re-indented.
% \leftskip indents plain paragraphs; \@totalleftmargin makes list environments
% (itemize/enumerate) start from the same indent instead of the page margin.
\makeatletter
\newlength{\lvlindent}
\setlength{\lvlindent}{1.8em}          % one indent step
\let\ulmsSecIndent\section
\renewcommand{\section}{\par\global\leftskip=0pt\global\@totalleftmargin=0pt\relax\ulmsSecIndent}
\let\ulmsSubsecIndent\subsection
\renewcommand{\subsection}{\par\global\leftskip=\lvlindent\global\@totalleftmargin=\lvlindent\relax\ulmsSubsecIndent}
\let\ulmsSubsubsecIndent\subsubsection
\renewcommand{\subsubsection}{\par\global\leftskip=2\lvlindent\global\@totalleftmargin=2\lvlindent\relax\ulmsSubsubsecIndent}
\makeatother

% Indent the headings themselves to match (titlesec ignores \leftskip for the
% heading line, so its left margin is set here via the first \titlespacing arg).
\titlespacing*{\subsection}   {\lvlindent} {1.0\baselineskip}{0.4\baselineskip}
\titlespacing*{\subsubsection}{2\lvlindent}{0.8\baselineskip}{0.3\baselineskip}

% ---- Lists (tighter, custom bullets) ----------------------------------------
\usepackage{enumitem}
\setlist[itemize]{leftmargin=1.6em, itemsep=2pt, topsep=3pt}
\setlist[enumerate]{leftmargin=1.9em, itemsep=2pt, topsep=3pt}

% ---- Tables ------------------------------------------------------------------
\usepackage{array}
\usepackage{booktabs}         % \toprule / \midrule / \bottomrule
\usepackage{tabularx}         % auto-width columns that wrap long Thai text
\usepackage{longtable}        % tables that can break across pages
\usepackage{multirow}         % merged cells for the team-structure table

% Left-aligned wrapping column types:
%   L{width} — fixed-width ragged-right paragraph column
%   Y        — flexible ragged-right X column (for tabularx)
\newcolumntype{L}[1]{>{\raggedright\arraybackslash}p{#1}}
\newcolumntype{Y}{>{\raggedright\arraybackslash}X}
\renewcommand{\arraystretch}{1.3}

% Helpers for a styled header row inside tables.
%   \hrow  — starts a header row and shades it (must be FIRST in the row;
%            colortbl allows only one \rowcolor per row).
%   \thead — bold header-cell text.
\newcommand{\hrow}{\rowcolor{headbg}}
\newcommand{\thead}[1]{\textbf{#1}}

% \begin{keeptable} … \end{keeptable} — keeps a table and the bold label above
% it on the SAME page (used by the data dictionary in §5). A tabularx cannot
% split across pages, so LaTeX is free to break between the label and its
% table, stranding the label alone at the bottom of a page; wrapping the pair
% in a minipage makes them one unbreakable block.
%   * \leftskip is zeroed inside: the minipage box is already placed at the
%     current heading indent, so leaving it set would indent the content twice
%     and push the table past the right margin.
%   * \parskip is restored because minipage's \@parboxrestore zeroes it, which
%     would glue the table straight onto the label line.
% Tables inside must use \linewidth (= the minipage width), not \textwidth.
\newenvironment{keeptable}
  {\par\medskip\noindent
   \begin{minipage}{\dimexpr\textwidth-\leftskip\relax}%
   \leftskip=0pt\relax
   \setlength{\parskip}{0.5\baselineskip plus 2pt}}
  {\end{minipage}\par\medskip}

% ---- Table of contents styling ----------------------------------------------
% By default LaTeX gives dotted leaders to subsections but NOT to sections
% (they print bold with just a right-aligned page number). tocloft lets us turn
% on dots for the section level too, so every TOC line runs "….." to its page.
\usepackage{tocloft}
\renewcommand{\cftsecdotsep}{\cftdotsep}   % dotted leaders for section entries
\renewcommand{\cftsecleader}{\cftdotfill{\cftsecdotsep}}
\renewcommand{\cftsecpagefont}{\normalfont}  % keep page numbers un-bold

% ---- Figures -----------------------------------------------------------------
\usepackage{graphicx}
\graphicspath{{figures/}}

% Landscape (rotated) full-page floats — needed by the ER diagram in §5, which
% is far wider than it is tall and would be unreadable scaled to \textwidth.
\usepackage{rotating}

% ---- Flowcharts / diagrams (TikZ) -------------------------------------------
% Used for the borrow–return workflow diagram (§5). Libraries: geometric shapes
% for decision diamonds, arrows.meta for arrowheads, positioning for relative
% node placement, calc/fit/backgrounds for routing loops and grouping.
\usepackage{tikz}
\usetikzlibrary{shapes.geometric, arrows.meta, positioning, calc, fit, backgrounds}

% ---- Links & bookmarks -------------------------------------------------------
\usepackage{hyperref}
\hypersetup{
  colorlinks=true, linkcolor=black, citecolor=linkblue, urlcolor=linkblue,
  pdflang={th},
}

% Thai labels for auto-generated headings.
\renewcommand{\contentsname}{สารบัญ (Table of Contents)}
\renewcommand{\tablename}{ตาราง}
\renewcommand{\figurename}{รูปที่}

:
%    1) \hypersetup{pdftitle={...}}  — ชื่อเอกสารใน metadata ของ PDF
%    2) \begin{document}
%  (pdftitle ถูกย้ายออกจากไฟล์นี้ เพราะเป็นค่าเฉพาะของแต่ละเอกสาร)
%
%  แก้ที่นี่ที่เดียว = มีผลกับทุกเอกสาร ห้ามก๊อป preamble ไปวางซ้ำ
% =============================================================================
% =============================================================================
%  main.tex  —  เอกสารข้อเสนอโครงการ (Academic Project Proposal)
%  ระบบบริหารจัดการการยืม–คืนอุปกรณ์มหาวิทยาลัย (ULMs)
%
%  IMPORTANT: This document contains Thai text and MUST be compiled with
%  XeLaTeX (not pdflatex). Thai needs a Unicode font + ICU line-breaking.
%      Build:  latexmk main.tex        (uses .latexmkrc -> xelatex)
%      or:     xelatex main.tex        (run twice for the table of contents)
% =============================================================================

\documentclass[11pt, a4paper]{article}

% ---- Fonts & Thai support ----------------------------------------------------
% fontspec lets XeLaTeX use any system font. Laksaman is a complete serif face
% that covers BOTH Thai and Latin glyphs, so one \setmainfont handles the whole
% (Thai + English) document without font switching.
\usepackage{fontspec}

% Thai has no spaces between words, so TeX cannot find line-break points on its
% own. These two XeTeX primitives switch on ICU's dictionary-based Thai word
% segmentation, which is what makes Thai paragraphs wrap correctly.
\XeTeXlinebreaklocale "th"
\XeTeXlinebreakskip = 0pt plus 0.1pt

\setmainfont{Laksaman}

% ---- Page geometry & paragraph style ----------------------------------------
\usepackage[margin=2.5cm]{geometry}
\usepackage{parskip}          % block paragraphs (no indent, gap between) — common
                              % and more readable for Thai documents
\linespread{1.15}             % a little extra leading; Thai has tall tone marks

% ---- Colours (monochrome / grayscale palette) ------------------------------
\usepackage[table]{xcolor}    % 'table' option enables \rowcolor for table headers
% Monochrome (grayscale) palette. Names kept as-is so every reference still
% works; only the tones changed from the original green/blue accents.
\definecolor{kugreen}{HTML}{1A1A1A}   % near-black — section headings & title
\definecolor{kudark}{HTML}{333333}    % dark gray — subsection/subsubsection titles
\definecolor{headbg}{HTML}{EAEAEA}    % light gray background for table header rows
\definecolor{linkblue}{HTML}{404040}  % medium-dark gray — links / citations

% ---- Section heading styling -------------------------------------------------
% 'nobottomtitles*' stops a heading from being printed near the bottom of a
% page: if less than \bottomtitlespace of room is left, the heading (and its
% following text) moves to the next page. This is titlesec's built-in, robust
% replacement for the earlier \needspace hack — the latter injected vertical
% glue that clashed with longtable page-splitting ("Infinite glue shrinkage").
\usepackage[nobottomtitles*]{titlesec}
\titleformat{\section}
  {\color{kugreen}\Large\bfseries}{\thesection}{0.6em}{}
\titleformat{\subsection}
  {\color{kudark}\large\bfseries}{\thesubsection}{0.5em}{}
\titleformat{\subsubsection}
  {\color{kudark}\normalsize\bfseries}{\thesubsubsection}{0.5em}{}
\titlespacing*{\section}{0pt}{1.4\baselineskip}{0.5\baselineskip}

% Minimum room a heading needs at the bottom of a page; if only ~1–3 lines are
% left (less than this), the heading is bumped to the next page instead of being
% stranded above its table/text.
\renewcommand{\bottomtitlespace}{4\baselineskip}

% ---- Hierarchical indentation by heading level ------------------------------
% Each heading level shifts BOTH the heading and the body text under it further
% to the right: section = 0, subsection (x.y) = 1 step, subsubsection (x.y.z) =
% 2 steps. We do this by having the sectioning commands set \leftskip (a global
% left margin that carries through every following line until the next heading
% resets it). \par first, so the *previous* paragraph is not re-indented.
% \leftskip indents plain paragraphs; \@totalleftmargin makes list environments
% (itemize/enumerate) start from the same indent instead of the page margin.
\makeatletter
\newlength{\lvlindent}
\setlength{\lvlindent}{1.8em}          % one indent step
\let\ulmsSecIndent\section
\renewcommand{\section}{\par\global\leftskip=0pt\global\@totalleftmargin=0pt\relax\ulmsSecIndent}
\let\ulmsSubsecIndent\subsection
\renewcommand{\subsection}{\par\global\leftskip=\lvlindent\global\@totalleftmargin=\lvlindent\relax\ulmsSubsecIndent}
\let\ulmsSubsubsecIndent\subsubsection
\renewcommand{\subsubsection}{\par\global\leftskip=2\lvlindent\global\@totalleftmargin=2\lvlindent\relax\ulmsSubsubsecIndent}
\makeatother

% Indent the headings themselves to match (titlesec ignores \leftskip for the
% heading line, so its left margin is set here via the first \titlespacing arg).
\titlespacing*{\subsection}   {\lvlindent} {1.0\baselineskip}{0.4\baselineskip}
\titlespacing*{\subsubsection}{2\lvlindent}{0.8\baselineskip}{0.3\baselineskip}

% ---- Lists (tighter, custom bullets) ----------------------------------------
\usepackage{enumitem}
\setlist[itemize]{leftmargin=1.6em, itemsep=2pt, topsep=3pt}
\setlist[enumerate]{leftmargin=1.9em, itemsep=2pt, topsep=3pt}

% ---- Tables ------------------------------------------------------------------
\usepackage{array}
\usepackage{booktabs}         % \toprule / \midrule / \bottomrule
\usepackage{tabularx}         % auto-width columns that wrap long Thai text
\usepackage{longtable}        % tables that can break across pages
\usepackage{multirow}         % merged cells for the team-structure table

% Left-aligned wrapping column types:
%   L{width} — fixed-width ragged-right paragraph column
%   Y        — flexible ragged-right X column (for tabularx)
\newcolumntype{L}[1]{>{\raggedright\arraybackslash}p{#1}}
\newcolumntype{Y}{>{\raggedright\arraybackslash}X}
\renewcommand{\arraystretch}{1.3}

% Helpers for a styled header row inside tables.
%   \hrow  — starts a header row and shades it (must be FIRST in the row;
%            colortbl allows only one \rowcolor per row).
%   \thead — bold header-cell text.
\newcommand{\hrow}{\rowcolor{headbg}}
\newcommand{\thead}[1]{\textbf{#1}}

% \begin{keeptable} … \end{keeptable} — keeps a table and the bold label above
% it on the SAME page (used by the data dictionary in §5). A tabularx cannot
% split across pages, so LaTeX is free to break between the label and its
% table, stranding the label alone at the bottom of a page; wrapping the pair
% in a minipage makes them one unbreakable block.
%   * \leftskip is zeroed inside: the minipage box is already placed at the
%     current heading indent, so leaving it set would indent the content twice
%     and push the table past the right margin.
%   * \parskip is restored because minipage's \@parboxrestore zeroes it, which
%     would glue the table straight onto the label line.
% Tables inside must use \linewidth (= the minipage width), not \textwidth.
\newenvironment{keeptable}
  {\par\medskip\noindent
   \begin{minipage}{\dimexpr\textwidth-\leftskip\relax}%
   \leftskip=0pt\relax
   \setlength{\parskip}{0.5\baselineskip plus 2pt}}
  {\end{minipage}\par\medskip}

% ---- Table of contents styling ----------------------------------------------
% By default LaTeX gives dotted leaders to subsections but NOT to sections
% (they print bold with just a right-aligned page number). tocloft lets us turn
% on dots for the section level too, so every TOC line runs "….." to its page.
\usepackage{tocloft}
\renewcommand{\cftsecdotsep}{\cftdotsep}   % dotted leaders for section entries
\renewcommand{\cftsecleader}{\cftdotfill{\cftsecdotsep}}
\renewcommand{\cftsecpagefont}{\normalfont}  % keep page numbers un-bold

% ---- Figures -----------------------------------------------------------------
\usepackage{graphicx}
\graphicspath{{figures/}}

% Landscape (rotated) full-page floats — needed by the ER diagram in §5, which
% is far wider than it is tall and would be unreadable scaled to \textwidth.
\usepackage{rotating}

% ---- Flowcharts / diagrams (TikZ) -------------------------------------------
% Used for the borrow–return workflow diagram (§5). Libraries: geometric shapes
% for decision diamonds, arrows.meta for arrowheads, positioning for relative
% node placement, calc/fit/backgrounds for routing loops and grouping.
\usepackage{tikz}
\usetikzlibrary{shapes.geometric, arrows.meta, positioning, calc, fit, backgrounds}

% ---- Links & bookmarks -------------------------------------------------------
\usepackage{hyperref}
\hypersetup{
  colorlinks=true, linkcolor=black, citecolor=linkblue, urlcolor=linkblue,
  pdflang={th},
}

% Thai labels for auto-generated headings.
\renewcommand{\contentsname}{สารบัญ (Table of Contents)}
\renewcommand{\tablename}{ตาราง}
\renewcommand{\figurename}{รูปที่}



\hypersetup{pdftitle={รายงานความก้าวหน้าโครงการ ULMs ครั้งที่ 3}}

% ---- ย่อหน้า (เหมือนฉบับที่ 2) ------------------------------------------------
\setlength{\parindent}{1.5em}
\setlength{\parskip}{0.35\baselineskip plus 2pt}
\titlespacing{\section} {0pt} {1.4\baselineskip}{0.5\baselineskip}
\titlespacing{\subsection} {\lvlindent} {1.0\baselineskip}{0.4\baselineskip}
\titlespacing{\subsubsection}{2\lvlindent}{0.8\baselineskip}{0.3\baselineskip}

\setlist[itemize]{leftmargin=1.6em, itemsep=2pt, topsep=0.55\baselineskip}
\setlist[enumerate]{leftmargin=1.9em, itemsep=2pt, topsep=0.55\baselineskip}

\setlength{\emergencystretch}{3em}
\newcommand{\zb}{\hspace{0pt}}

\usepackage{float}

% ---- ตัวแปรของรอบรายงาน ------------------------------------------------------
\newcommand{\reportdate}{10 กันยายน 2569}
\newcommand{\reportperiod}{4 – 10 กันยายน 2569}
\newcommand{\actualpct}{75\%}
\newcommand{\planpct}{78\%}
\newcommand{\prevactualpct}{70\%}

\begin{document}

% =============================================================================
% TITLE PAGE
% =============================================================================
\begin{titlepage}
\centering
\vspace*{2.8cm}

{\Large\bfseries\color{kugreen} รายงานความก้าวหน้าโครงการ ครั้งที่ 3\par}
\vspace{0.6cm}
{\LARGE\bfseries\color{kugreen} ระบบบริหารจัดการการยืม–คืนอุปกรณ์มหาวิทยาลัย\par}
\vspace{0.35cm}
{\Large\color{kudark} University Lending Management System (ULMs)\par}

\vspace{1.8cm}
\rule{0.72\textwidth}{0.8pt}
\vspace{0.9cm}

\begin{tabular}{@{}r@{\hspace{1.2em}}l@{}}
\textbf{รอบรายงาน} & \reportperiod \\[0.4em]
\textbf{วันที่รายงาน} & \reportdate \\[0.4em]
\textbf{ขอบเขตการนับ} & งานพัฒนาระบบ Sprint 0–5 (ไม่รวมงานก่อน 24 ก.ค.) \\[0.4em]
\textbf{สถานะ Sprint} & วันสิ้นสุด Sprint 3 (Sprint ที่ 4 จาก 6) \\[0.4em]
\textbf{ความก้าวหน้า} & \actualpct\ เทียบแผน \planpct\ \ (ต่ำกว่าแผน 3\%) \\[0.4em]
\textbf{งบประมาณ} & ใช้ไป 0 บาท จากวงเงินประเมิน 5{,}000 บาท \\[0.4em]
\textbf{กำหนดเสร็จ} & 1 ตุลาคม 2569 (ยังไม่มีการเลื่อน) \\
\end{tabular}

\vspace{0.9cm}
\rule{0.72\textwidth}{0.8pt}

\vfill
{\large\color{kudark} คณะทำงาน ``miro ปั่น 14 บาท''\par}
\vspace{0.3cm}
{\color{kudark} ภาควิชาวิศวกรรมคอมพิวเตอร์ คณะวิศวกรรมศาสตร์\par}
{\color{kudark} มหาวิทยาลัยเกษตรศาสตร์\par}
\vspace{1.2cm}
\end{titlepage}

% =============================================================================
\section{บทสรุป}

รอบรายงานนี้คือสัปดาห์ที่สองและเป็นวันสิ้นสุดของ Sprint 3 ความก้าวหน้าเพิ่มจาก \prevactualpct\ เป็น \actualpct\ ขณะที่แผนขยับจาก 72\% เป็น \planpct\ ระยะห่างจากแผนจึงกว้างขึ้นเล็กน้อยจาก 2\% เป็น 3\%

Sprint 3 ไม่ผ่านเกณฑ์ ``เสร็จ'' ที่ตั้งไว้ คือการเดินวงจรยืม–คืนครบเส้นทางปกติผ่านหน้าจอจริง จุดวัดสามข้อที่มีกำหนดภายในรอบนี้ไม่ผ่านทั้งสามข้อ สาเหตุเดียวกันทั้งหมดคือปุ่มส่งคำขอของผู้ยืมยังไม่ได้ต่อเข้ากับ backend และหน้าส่งมอบของเจ้าหน้าที่ยังเป็นหน้าว่าง ทั้งที่ procedure ปลายทางมีอยู่บน main ครบตั้งแต่ 2 กันยายน

ความก้าวหน้าที่เกิดขึ้นจริงในรอบนี้อยู่นอกเส้นทางนั้นเกือบทั้งหมด ฝั่ง backend เพิ่มระบบต่ออายุการยืมครบทั้งกลุ่มความต้องการ พร้อมตารางเลือกผู้ตัดสินที่มี test ของตัวเอง cron job ทำงานตามเวลาได้ 5 จาก 8 งาน มีรายงานสรุปรายหน่วยงาน และบังคับกฎ ``ผู้ตรวจสภาพ T2 ต้องไม่ใช่ผู้จัดเตรียม'' ที่ชั้น service แล้ว ฝั่ง frontend มีหน้าจอที่อ่านข้อมูลจริงเพิ่มจาก 13 เป็น 21 หน้า และคำแปลภาษาไทยของข้อความผิดพลาดครอบคลุมเพิ่มจาก 29 เป็น 60 จาก 63 รหัส

การตรวจสอบเพื่อทำรายงานนี้พบข้อถดถอยหนึ่งข้อ นับจาก 6 กันยายน หน้า ``คำขอของฉัน'' อ่านรายการจาก backend แทน store ในเบราว์เซอร์ แต่ปุ่มส่งคำขอยังเขียนลง store อยู่ คำขอยืมอุปกรณ์ที่ผู้ยืมเพิ่งกดส่งจึงไม่ปรากฏในรายการเลย ก่อนหน้านี้อย่างน้อยยังเห็นได้จนกว่าจะรีเฟรช ข้อนี้ทำให้งานต่อปุ่มส่งคำขอเป็นงานลำดับแรกของ Sprint 4

% =============================================================================
\section{แผนภาพรวมทั้ง 6 Sprint}
\label{sec:overview}

\noindent\begin{tabularx}{\dimexpr\textwidth-\leftskip\relax}{@{} L{2.9cm} Y L{2.2cm} @{}}
\toprule
\hrow \thead{Sprint} & \thead{ธีมและงานหลัก} & \thead{สถานะ} \\
\midrule
\textbf{Sprint 0} \newline 24 – 30 ก.ค. & วางฐาน — ปิดขอบเขต · ติดตั้งโครงสร้างพื้นฐานทั้งสองฝั่ง · database schema แกนกลาง · ระบบออกแบบและ mock & ผ่าน \\
\textbf{Sprint 1} \newline 31 ก.ค. – 13 ส.ค. & ล็อก contract และเปิดทาง — ล็อก contract API · ระบบเข้าสู่ระบบและออกจากระบบ · CI · เกณฑ์การยอมรับของวงจรหลัก & ผ่าน \\
\textbf{Sprint 2} \newline 14 – 27 ส.ค. & วงจรยืม–คืนของอุปกรณ์เคลื่อนย้ายได้ — รายการอุปกรณ์และจำนวนคงเหลือ · สร้างคำขอ อนุมัติ ส่งมอบ รับคืน · ข้อมูล seed · เลิกใช้ mock ในเส้นทางหลัก & ไม่ผ่านเกณฑ์ \newline (ยกงานไป Sprint 3) \\
\textbf{Sprint 3} \newline 28 ส.ค. – 10 ก.ย. & กฎธุรกิจรายระดับและการจองสถานที่ (T3) — อนุมัติตามระดับอุปกรณ์ · ผูกหมายเลขชิ้น · สล็อตเวลาและการจองห้อง · ตรวจสภาพและระบบเครดิต \newline รับงานที่ยกมาจาก Sprint 2 เป็นลำดับแรก & ไม่ผ่านเกณฑ์ \newline (กฎธุรกิจครบฝั่ง backend แต่วงจรยังเดินผ่านหน้าจอไม่ได้ ยกงาน 4 วันไป Sprint 4) \\
\textbf{Sprint 4} \newline 11 – 24 ก.ย. & ข้ามหน่วยงานและงานขัดเงา — ตะกร้าข้ามหน่วยงาน · จองล่วงหน้า · ต่ออายุ · หน้าจอที่เหลือ · การค้นหาและการรองรับมือถือ \newline สี่วันแรกรับงานค้างของ Sprint 3 ก่อน & เริ่มพรุ่งนี้ \newline (ต่ออายุฝั่ง backend ทำล่วงหน้าเสร็จแล้ว) \\
\textbf{Sprint 5} \newline 25 ก.ย. – 1 ต.ค. & หยุดเพิ่มความสามารถ — ทดสอบถดถอย · นำขึ้นใช้งานจริง · คู่มือผู้ใช้และคู่มือติดตั้ง · ซ้อมนำเสนอ & ตามแผน \\
\bottomrule
\end{tabularx}

% =============================================================================
\section{สถานะโครงการโดยรวม}

\noindent\begin{tabularx}{\dimexpr\textwidth-\leftskip\relax}{@{} L{1.5cm} L{3.0cm} L{2.4cm} Y @{}}
\toprule
\hrow \thead{ด้าน} & \thead{แผน} & \thead{ผลจริง} & \thead{การประเมิน} \\
\midrule
เวลา & ผ่านไป 6.9 จาก 10 สัปดาห์ (69\%) & ผ่านไป 6.9 จาก 10 สัปดาห์ (69\%) & ตามปฏิทิน วันเสร็จยังเป็น 1 ตุลาคม 2569 แต่ Sprint 3 ปิดโดยมีงานค้าง 4 วันทำงานยกไป Sprint 4 \\
งาน & \planpct & \actualpct & ต่ำกว่าแผน 3\% กว้างขึ้นจาก 2\% ในรอบก่อน งานในรอบนี้เดินหน้าที่ backend และหน้าจอฝั่งผู้ดูแลระบบ แต่มิติ Test แทบไม่ขยับ รายละเอียดอยู่ในหัวข้อ \ref{sec:dimension} \\
เงิน & 5{,}000 บาท & 0 บาท (0\%) & ยังไม่มีรายจ่าย รายการที่ประเมินไว้คือค่าเช่า Cloud/Hosting และฐานข้อมูลกลาง ซึ่งจะเกิดขึ้นเมื่อนำระบบขึ้นใช้งานจริงใน Sprint 5 \\
\bottomrule
\end{tabularx}

% =============================================================================
\clearpage
\section{ความก้าวหน้าแยกตามมิติของงาน}
\label{sec:dimension}

ค่าในคอลัมน์ ``จริง'' ประเมินจากผลงานที่ตรวจสอบได้ใน repository และนับเฉพาะงานที่รวมเข้า main แล้ว วันรายงานรอบนี้ตรงกับวันสิ้นสุด Sprint 3 ค่าในคอลัมน์ ``แผน'' จึงอ่านตรงจากแถว S3 ของตารางภาคผนวก

\noindent\begin{tabularx}{\dimexpr\textwidth-\leftskip\relax}{@{} L{3.2cm} L{1.1cm} L{1.0cm} L{1.0cm} L{1.2cm} Y @{}}
\toprule
\hrow \thead{มิติงาน} & \thead{น้ำหนัก} & \thead{แผน} & \thead{จริง} & \thead{ผลต่าง} & \thead{ข้อเท็จจริงประกอบ} \\
\midrule
Back-End Design \newline (API + Database Model) & 10\% & 100\% & 97\% & $-3$ & contract ขยายเป็น 10 domain 86 procedure ออกแบบการต่ออายุครบพร้อมตารางเลือกผู้ตัดสิน และตาราง cron run log ยังเหลือการออกแบบฝั่ง backend ของการอุทธรณ์และการจองสล็อตเวลา \\
Front-End Design \newline (UI/UX + Business Process) & 10\% & 92\% & 95\% & $+3$ & หน้าจอที่มีเนื้อหาจริง 26 จาก 28 หน้า รวมหน้าจัดการผู้ใช้และสิทธิ์ของหน่วยงาน ยุบหน้าที่ซ้ำกัน 2 หน้า เหลือหน้าว่าง 2 หน้า คือส่งมอบของและส่งออกรายงาน \\
Back-End Dev & 30\% & 78\% & 77\% & $-1$ & เพิ่มการต่ออายุ 7 procedure cron job ทำงานตามเวลา 5 จาก 8 งาน รายงานสรุป และกฎผู้ตรวจต้องไม่ใช่ผู้จัดเตรียม ยังเหลือการอุทธรณ์ การจองสล็อต การแก้ค่าตั้งเชิงเทคนิค และ cron job อีก 3 งาน \\
Front-End Dev & 25\% & 78\% & 74\% & $-4$ & หน้าจอที่อ่านข้อมูลจริงเพิ่มจาก 13 เป็น 21 หน้า แต่ปุ่มส่งคำขอและปุ่มขอต่ออายุของผู้ยืมยังไม่เรียก backend และมีข้อถดถอยที่คำขอที่เพิ่งส่งหายจากรายการ ยังใช้ไฟล์ contract ที่เขียนเอง \\
Test (Front-End + Back-End) & 15\% & 52\% & 35\% & $-17$ & test ฝั่ง backend 234 กรณี เพิ่ม 7 กรณีจากตารางเลือกผู้ตัดสินการต่ออายุ งานตรวจเวอร์ชันใน CI ขวางการรวมโค้ดแล้ว แต่ยังไม่มี test ของวงจรยืม–คืนเลย service ที่มี test คือ 8 จาก 21 และฝั่ง frontend ยังไม่มี test runner \\
Documentation & 10\% & 76\% & 88\% & $+12$ & SDS ปรับเป็นฉบับ 1.1 ตามโค้ดของวันรายงาน รวมข้อต่างระหว่างโค้ดกับ SRS ที่ต้องตัดสิน และหลักฐานการเรียก API จริงของวงจรยืม–คืน \\
\midrule
\hrow \thead{รวมถ่วงน้ำหนัก} & \thead{100\%} & \thead{\planpct} & \thead{\actualpct} & \thead{$-3$} & \\
\bottomrule
\end{tabularx}

มิติที่ต้องเฝ้าระวังในรอบก่อนคือ Test และมิตินี้ห่างจากแผนมากขึ้นจาก $-13$ เป็น $-17$ แผนรอบนี้ขอให้ Test ขยับ 7 จุดจากงานสามรายการ คือ test ที่เดินวงจรยืม–คืนบนฐานข้อมูลจริง test runner ฝั่ง frontend และการย้ายชุด AC check เข้า CI ทั้งสามรายการไม่เกิดขึ้น สิ่งที่ขยับได้ 3 จุดมาจาก test ของตารางเลือกผู้ตัดสินการต่ออายุ และงานตรวจเวอร์ชันที่เปลี่ยนเป็นงานบังคับ

สัดส่วนของ service ที่มี test ลดลงจาก 8 ใน 18 เหลือ 8 ใน 21 เพราะ service ที่เพิ่มเข้ามาใหม่ในรอบนี้สามตัว คือการต่ออายุ cron job และรายงาน ไม่มี test ที่ระดับ service เลย ส่วนที่มี test คือกฎการเลือกผู้ตัดสินซึ่งแยกออกมาเป็นฟังก์ชันล้วน ไม่ใช่ตัว service ที่อ่านและเขียนฐานข้อมูล


\begin{figure}[H]
\centering
\begin{tikzpicture}[x=0.062cm, y=1cm]
 \foreach \x in {0,20,40,60,80,100} {
 \draw[gray!25] (\x,-0.30) -- (\x,6.15);
 \node[below, font=\scriptsize] at (\x,-0.30) {\x};
 }
 \node[below, font=\scriptsize] at (50,-0.68) {\% ความก้าวหน้าของมิตินั้น};
 \draw[thick] (0,-0.30) -- (0,6.15);

 \foreach \y/\name/\wt/\plan/\act in {%
 5.6/{Back-End Design}/10/100/97,
 4.6/{Front-End Design}/10/92/95,
 3.6/{Back-End Dev}/30/78/77,
 2.6/{Front-End Dev}/25/78/74,
 1.6/{Test}/15/52/35,
 0.6/{Documentation}/10/76/88} {
 \node[left, xshift=-0.16cm, font=\scriptsize, align=right] at (0,\y)
 {\name\\[-3pt]{\tiny น้ำหนัก \wt\%}};
 \fill[gray!45] (0,{\y+0.05}) rectangle (\plan,{\y+0.30});
 \fill[kugreen] (0,{\y-0.30}) rectangle (\act,{\y-0.05});
 \node[right, font=\tiny, gray!45!black] at (\plan,{\y+0.175}) {\plan};
 \node[right, font=\tiny, kugreen] at (\act,{\y-0.175}) {\act};
 }

 \fill[gray!45] (0,-1.36) rectangle (7,-1.54);
 \node[right, font=\scriptsize] at (9,-1.45) {แผน};
 \fill[kugreen] (0,-1.76) rectangle (7,-1.94);
 \node[right, font=\scriptsize] at (9,-1.85) {ผลจริง};
\end{tikzpicture}
\vspace{0.6em}
\caption{ความก้าวหน้ารายมิติ ณ \reportdate\ วันสิ้นสุด Sprint 3 เทียบกับค่าเป้าหมายตามแผน
 ตัวเลขใต้ชื่อมิติคือน้ำหนักที่ใช้ถ่วงเมื่อรวมเป็นค่าเดียว
 แท่งผลจริงนับเฉพาะงานที่รวมเข้า main แล้ว
 มิติที่ห่างจากแผนมากที่สุดคือ Test ที่ $-17$
 ส่วน Documentation กับ Front-End Design เดินนำแผน ที่ $+12$ และ $+3$}
\label{fig:dims}
\end{figure}

% =============================================================================
\clearpage
\section{กราฟความก้าวหน้าเทียบแผน}

เส้นแผนสร้างจากแผน Sprint โดยตรง คือกำหนดค่าเป้าหมายรายมิติ ณ วันสิ้นสุดของแต่ละ Sprint แล้วคูณน้ำหนักรวมกันเป็นค่าสะสม ตารางค่าทั้งหมดอยู่ในภาคผนวก แกนนอนเริ่มที่ 24 กรกฎาคม 2569 ซึ่งเป็นวันเริ่มลงมือพัฒนา

ความชันของเส้นผลจริงในรอบนี้ลดลงจากรอบก่อน จากรอบก่อนขยับ 18 จุดในแปดวัน ไปรอบนี้ขยับ 5 จุดในหกวัน ซึ่งตรงกับที่คาดไว้ว่าความชันของรอบก่อนนั้นมีส่วนหนึ่งมาจากงานที่เขียนเสร็จก่อนแล้วเพิ่งถูกนับ

\begin{figure}[H]
\centering
\begin{tikzpicture}[x=1.20cm, y=0.058cm]
 \foreach \y in {0,20,40,60,80,100}
 \draw[gray!22] (0,\y) -- (10.2,\y);

 \foreach \x in {0.86, 2.86, 4.86, 6.86, 8.86}
 \draw[gray!40, dashed, thin] (\x,0) -- (\x,108);

 \foreach \x/\lbl in {0.43/{S0}, 1.86/{S1}, 3.86/{S2}, 5.86/{S3}, 7.86/{S4}, 9.36/{S5}}
 \node[font=\scriptsize\bfseries, gray!50!black] at (\x,114) {\lbl};

 \draw[->, thick] (0,0) -- (10.6,0);
 \draw[->, thick] (0,0) -- (0,122) node[above, font=\scriptsize] {\% ความก้าวหน้าสะสม};
 \foreach \y in {0,20,40,60,80,100}
 \draw (0,\y) -- (-0.10,\y) node[left, font=\scriptsize] {\y};

 \foreach \x/\lbl in {0/{24 ก.ค.}, 0.86/{30 ก.ค.}, 2.86/{13 ส.ค.},
 4.86/{27 ส.ค.}, 6.86/{10 ก.ย.}, 8.86/{24 ก.ย.}, 9.86/{1 ต.ค.}}
 {\draw (\x,0) -- (\x,-3);
 \node[rotate=45, anchor=north east, font=\scriptsize] at (\x,-4) {\lbl};}

 % ---- เส้นแผน ----
 \draw[thick, kudark, densely dashed]
 plot coordinates {(0,1) (0.86,12) (2.86,42) (4.86,65) (6.86,78) (8.86,94) (9.86,100)};
 \foreach \p in {(0,1), (0.86,12), (2.86,42), (4.86,65), (6.86,78), (8.86,94), (9.86,100)}
 \fill[kudark] \p circle (1.6pt);

 % ---- เส้นผลจริง — จุดคือวันที่วัด 6.86 = 10 ก.ย. = วันรายงานนี้ ----
 \draw[very thick, kugreen]
 plot coordinates {(0,0) (2.0,24) (2.43,32) (2.86,43) (4.29,43) (4.86,52) (6.0,70) (6.86,75)};
 \foreach \p in {(0,0), (2.0,24), (2.43,32), (2.86,43), (4.29,43), (4.86,52), (6.0,70), (6.86,75)}
 \fill[kugreen] \p circle (2.2pt);

 \draw[gray!70, thick, |-|] (6.86,75) -- (6.86,78);
 \node[right, font=\tiny, gray!70!black] at (6.92,76.5) {$-3$};

 % sits above both lines so its white box hides neither
 \node[anchor=south east, font=\scriptsize, gray!55!black, align=right,
 fill=white, inner sep=1.5pt] at (5.95,82)
 {4 ก.ย.\\70\% · แผน 72\%};
 \node[anchor=north west, font=\scriptsize, kugreen, align=left,
 fill=white, inner sep=1.5pt] at (6.95,70)
 {10 ก.ย. (วันรายงาน)\\ผลจริง 75\% · แผน 78\%};

 \fill[white] (5.85,4) rectangle (10.2,26);
 \draw[thick, kudark, densely dashed] (6.0,19) -- (6.7,19);
 \node[right, font=\scriptsize] at (6.8,19) {แผน (Planned)};
 \draw[very thick, kugreen] (6.0,9) -- (6.7,9);
 \node[right, font=\scriptsize] at (6.8,9) {ผลจริง (Actual)};
\end{tikzpicture}
\caption{กราฟ เปรียบเทียบความก้าวหน้าตามแผนกับผลการดำเนินงานจริง
 แกนนอนวางตามขอบ Sprint ทั้ง 6 ช่วง (เส้นประแนวตั้ง)
 จุดบนเส้นผลจริงคือวันที่วัดความก้าวหน้า
 วันรายงาน 10 กันยายน 2569 ตรงกับขอบ Sprint 3 พอดี
 ช่วงราบระหว่าง 13 ถึง 23 สิงหาคม คือช่วงสอบกลางภาค}
\label{fig:scurve}
\end{figure}

% =============================================================================
% the 9-row table cannot break, so start the section on a fresh page rather
% than leave its heading stranded under the S-curve
\clearpage
\section{ผลการดำเนินงานเทียบแผนรายวันที่ตั้งไว้}
\label{sec:done}

รายงานความก้าวหน้าครั้งที่ 2 กำหนดงานรายวันไว้ 9 รายการสำหรับช่วงที่เหลือของ Sprint 3 ตารางนี้รายงานผลของทั้ง 9 รายการตามลำดับเดิม

\noindent\begin{tabularx}{\dimexpr\textwidth-\leftskip\relax}{@{} L{1.3cm} Y L{1.9cm} L{3.4cm} @{}}
\toprule
\hrow \thead{กำหนด} & \thead{งานที่ตั้งไว้} & \thead{ผู้รับผิดชอบ} & \thead{ผล} \\
\midrule
4 ก.ย. & สร้างไฟล์ contract จากโค้ดจริงเข้า main และทำให้ฝั่ง frontend อ้างถึง schema ได้ & BE-lead & ยังไม่ได้ทำ ตัวสร้างยังทำงานได้ แต่ไฟล์ที่สร้างยังถูกตั้งให้ git ไม่เก็บ \\
5 ก.ย. & รับไฟล์ contract ไปใช้แทนไฟล์ที่ฝั่ง frontend เขียนเอง & FE-lead & ยังไม่ได้ทำ แต่ไฟล์ที่เขียนเองถูกแก้ให้ตรงขึ้น procedure ที่เรียกไม่ได้ลดจาก 15 เหลือ 11 \\
6 ก.ย. & ต่อปุ่มส่งคำขอของผู้ยืมเข้ากับ procedure สร้างคำขอจริง & FE-dev & ยังไม่ได้ทำ ต่อเฉพาะการอ่านรายการและการยกเลิก ซึ่งทำให้เกิดข้อถดถอยตามบทสรุป \\
7 ก.ย. & หน้าส่งมอบของเจ้าหน้าที่ & FE-lead & ยังไม่ได้ทำ ยังเป็นหน้าว่าง \\
8 ก.ย. & test ที่เดินวงจรยืม–คืนบนฐานข้อมูลจริง & BE-lead & ยังไม่ได้ทำ \\
8 ก.ย. & ติดตั้ง test runner ฝั่ง frontend ย้าย AC check เข้า CI และเปลี่ยนงานตรวจเวอร์ชันให้ขวางการรวมโค้ด & QA-lead & ทำได้หนึ่งในสาม งานตรวจเวอร์ชันขวางการรวมโค้ดแล้ว เมื่อ 10 ก.ย. \\
9 ก.ย. & ต่อหน้ารายการยืมของฉัน หน้ารายการห้อง และหน้ารับของ เข้ากับข้อมูลจริง & FE-lead & ทำได้สองในสาม รายการยืมของฉันและหน้ารับของอ่านจริงแล้วตั้งแต่ 6 ก.ย. หน้าห้องยังเป็น mock \\
9 ก.ย. & กลุ่มการอุทธรณ์ฝั่ง backend และที่เก็บค่าตั้งเวลาทำการ & BE-dev & ยังไม่ได้ทำ ค่าตั้งเชิงเทคนิคอ่านได้แล้วแต่ยังแก้ไม่ได้ \\
10 ก.ย. & เดินวงจรยืม–คืนครบเส้นทางปกติผ่านหน้าจอจริง (เกณฑ์เสร็จของ Sprint 3) & PM & ไม่ผ่าน ติดที่ปุ่มส่งคำขอและหน้าส่งมอบ \\
\bottomrule
\end{tabularx}

สรุปผลของทั้ง 9 รายการคือเสร็จครบ 0 รายการ ทำได้บางส่วน 3 รายการ และยังไม่ได้ทำ 6 รายการ ผลนี้ต่างจากรอบก่อนที่เสร็จ 5 จาก 9 อย่างชัดเจน
\clearpage
ข้อสังเกตที่คณะทำงานบันทึกไว้คือ 7 จาก 9 รายการเป็นงานต่อสิ่งที่มีอยู่แล้วเข้าหากัน ไม่ใช่งานสร้างของใหม่ ขณะที่งานที่เกิดขึ้นจริงในรอบนี้ส่วนใหญ่เป็นงานสร้างของใหม่ที่อยู่ในแผน Sprint 4 งานต่อเชื่อมจึงถูกเลื่อนออกไป

\noindent\begin{tabularx}{\dimexpr\textwidth-\leftskip\relax}{@{} L{1.3cm} Y L{1.9cm} L{1.4cm} @{}}
\toprule
\hrow \thead{วันที่} & \thead{งานที่ทำเพิ่มนอกแผน} & \thead{ผู้รับผิดชอบ} & \thead{ผล} \\
\midrule
6 ก.ย. & หน้าจอผู้ดูแลระบบ 4 หน้า คือภาพรวม สถานะระบบ รายงาน และค่าตั้งเชิงเทคนิค เลิกใช้ mock ทั้งหมด พร้อม procedure รายงานสรุปรายหน่วยงานฝั่ง backend & FE-lead & เสร็จ \\
6 ก.ย. & หน้าจัดการผู้ใช้และสิทธิ์ของหน่วยงาน ซึ่งเป็นหน้าว่างในรอบก่อน & FE-lead & เสร็จ \\
6 ก.ย. & cron job ทำงานตามเวลาจริง 5 จาก 8 งาน พร้อมตารางบันทึกผลการรันทุกครั้ง และหน้าสั่งรันเองของผู้ดูแลระบบ & FE-lead & เสร็จ \\
6 ก.ย. & ตรึงเวอร์ชัน Node ของทั้งทีมไว้ในไฟล์เดียว และแก้ lockfile ให้ตรงกับเวอร์ชันที่กำหนด & FE-lead & เสร็จ \\
9 ก.ย. & ระบบต่ออายุการยืมฝั่ง backend ครบ 7 procedure พร้อมตารางเลือกผู้ตัดสินและ test 7 กรณี (งานของ Sprint 4) & BE-dev & เสร็จ \\
9 ก.ย. & บังคับกฎผู้ตรวจสภาพ T2 ต้องไม่ใช่ผู้จัดเตรียม (FR-RTN-04) & BE-dev & เสร็จ \\
9 ก.ย. & แก้คำแปลภาษาไทยของหน้าจอและข้อความผิดพลาด & FE-lead & เสร็จ \\
9 ก.ย. & ทดลองเรียก API จริง 33 ครั้งบนฐานข้อมูลเปล่า เดินวงจรยืม–คืนตั้งแต่ลงทะเบียนอุปกรณ์จนรับคืน & SA & เสร็จ \\
10 ก.ย. & ปรับ SDS เป็นฉบับ 1.1 ตามโค้ดของวันรายงาน & SA & เสร็จ \\
\bottomrule
\end{tabularx}

% =============================================================================
\section{ผลการทบทวนปิด Sprint 3}
\label{sec:review}

รายงานครั้งที่ 2 ประกาศจุดวัดไว้สี่ข้อ และระบุว่าถ้าข้อใดไม่ผ่านจะทบทวนลำดับความสำคัญของ Sprint 4 ทันที ตารางนี้คือผลของทั้งสี่ข้อ

\noindent\begin{tabularx}{\dimexpr\textwidth-\leftskip\relax}{@{} L{0.6cm} Y L{1.7cm} L{3.6cm} @{}}
\toprule
\hrow \thead{ข้อ} & \thead{จุดวัด} & \thead{กำหนด} & \thead{ผล} \\
\midrule
1 & ปุ่มส่งคำขอของผู้ยืมเขียนลงฐานข้อมูลจริง และหน้าส่งมอบของเจ้าหน้าที่ใช้งานได้ & 7 ก.ย. & ไม่ผ่านทั้งสองส่วน \\
2 & เดินวงจรยืม–คืนครบเส้นทางปกติผ่านหน้าจอจริง & 10 ก.ย. & ไม่ผ่าน เพราะข้อ 1 \\
3 & วงจรยืม–คืนมี test อัตโนมัติที่เดินครบเส้นทางบนฐานข้อมูลจริง & 10 ก.ย. & ไม่ผ่าน มีเพียงการเรียก API ด้วยมือหนึ่งรอบ \\
4 & ไม่เหลือ mock ในเส้นทางหลัก และฝั่ง frontend ใช้ไฟล์ contract ที่สร้างจากโค้ดจริง & 14 ก.ย. & ยังไม่ถึงกำหนด เหลือ mock 5 หน้า และยังไม่ได้เริ่มรับไฟล์ contract \\
\bottomrule
\end{tabularx}
\clearpage
เกณฑ์ที่ตั้งไว้ของ Sprint 3 เองคือกฎธุรกิจรายระดับกับการจองสถานที่ ฝั่ง backend ผ่านเกือบครบ การอนุมัติตามระดับอุปกรณ์ การผูกหมายเลขชิ้น การตรวจสภาพ และระบบเครดิตทำงานได้ และการเรียก API จริงเมื่อ 9 กันยายน ยืนยันว่าวงจรยืม–คืนเดินผ่าน HTTP ได้ครบตั้งแต่สร้างคำขอจนรับคืน ส่วนที่ไม่ผ่านคือทุกอย่างที่ต้องผ่านหน้าจอของผู้ยืม และการจองห้องซึ่งหน้าจอยังใช้ mock ทั้งสามหน้า

ผลการทบทวนมีสามข้อ ข้อแรก สี่วันทำงานแรกของ Sprint 4 คือ 11 ถึง 14 กันยายน ใช้กับงานค้างของ Sprint 3 เท่านั้น และไม่เริ่มความสามารถใหม่จนกว่าวงจรยืม–คืนจะเดินผ่านหน้าจอได้ ข้อที่สอง งานต่อเชื่อมแต่ละรายการมีผู้รับผิดชอบคนเดียวและรายงานผลในวันที่กำหนด แทนที่จะรวมไว้กับงานใหญ่ ข้อที่สาม การต่ออายุฝั่ง backend ที่ทำล่วงหน้าแล้วถูกนับเป็นงานของ Sprint 4 ที่เสร็จก่อนกำหนด เวลาที่ประหยัดได้ใช้ชดเชยสี่วันที่ยกมา

% =============================================================================
\section{ผลงานที่ส่งมอบในรอบนี้}

หัวข้อนี้รายงานเฉพาะสิ่งที่ทำงานได้จริงและตรวจสอบได้ใน repository

\subsection{ฝั่ง backend}

\begin{itemize}
 \item การต่ออายุการยืม — ผู้ยืมดูตัวเลือกก่อนยืนยัน ขอต่อ ดูคำขอของตนเอง และยกเลิกคำขอได้ · เจ้าหน้าที่และอาจารย์มีคิวพิจารณาคนละคิว · ระบบเลือกผู้ตัดสินจากระดับอุปกรณ์ ระดับเครดิต และครั้งที่ของการต่อ · กันช่วงเวลาที่ขอต่อไม่ให้ชนกับการจองของคนถัดไป · แจ้งผู้ยืมเมื่อได้รับอนุมัติ
 \item cron job ทำงานตามเวลาจริง — ทำเครื่องหมายคืนช้า แจ้งของหายเมื่อเลยกำหนดสองสัปดาห์ หมดอายุบทลงโทษ เตือนก่อนครบกำหนด และยกเลิกคำขอที่เลยกำหนดมารับ · ทุกครั้งที่รันบันทึกผลลงฐานข้อมูล และผู้ดูแลระบบสั่งรันเองได้
 \item รายงานสรุปรายหน่วยงาน — เจ้าหน้าที่เห็นเฉพาะหน่วยงานของตนเอง ผู้ดูแลระบบเห็นทั้งระบบ
 \item กฎแยกหน้าที่ของ T2 — ผู้ประเมินความเสียหายต้องไม่ใช่ผู้ที่จัดเตรียมชิ้นนั้น ระบบอ่านผู้จัดเตรียมจากบันทึกสภาพตอนจ่ายของ จึงไม่ต้องเพิ่มคอลัมน์ในฐานข้อมูล
 \item อาจารย์ทำงานของเจ้าหน้าที่ได้ทั้งหมด — แก้ middleware ที่เคยกันอาจารย์ออกจาก approval queue ทั้งที่ queue นั้นมีไว้ให้อาจารย์ใช้
 \item ขนาดของสิ่งที่ส่งมอบ ณ วันรายงาน — 10 domain รวม 86 procedure บนฐานข้อมูล 30 ตาราง 9 enum ผ่าน 9 migration
\end{itemize}

ข้อจำกัดที่เหลือมีสี่ข้อ ข้อแรกคือยังไม่มี test ของวงจรยืม–คืน การอนุมัติ การต่ออายุ และ notification ที่ระดับ service หลักฐานที่มีคือการเรียก API ด้วยมือหนึ่งรอบ ข้อที่สองคือการอุทธรณ์และการจองสล็อตเวลายังไม่มีฝั่ง backend ข้อที่สามคือการต่ออายุต่างจาก SRS สองจุด ผู้ยืมเครดิต D2 ถึง D3 ถูกส่งไปให้เจ้าหน้าที่ตรวจของ ขณะที่ FR-RNW-06 ให้อาจารย์อนุมัติ และห้อง T3 ถูกส่งให้เจ้าหน้าที่ทุกครั้ง ขณะที่ FR-RNW-04 ให้ต่อได้ทันทีเมื่อสล็อตถัดไปว่าง ต้องตัดสินว่าจะแก้โค้ดหรือแก้ SRS ข้อที่สี่คือไฟล์รูปหลักฐานยังเปิดอ่านได้โดยไม่ตรวจสิทธิ์ เหมือนรอบก่อน

\subsection{ฝั่ง frontend}

\begin{itemize}
 \item หน้าจอที่อ่านข้อมูลจริงเพิ่มเป็น 21 จาก 28 หน้า — เพิ่มหน้าผู้ดูแลระบบ 4 หน้า หน้าจัดการผู้ใช้และสิทธิ์ของหน่วยงาน 2 หน้า รายการยืมของฉัน และหน้ารับของ
 \item หน้าจอที่มีเนื้อหาจริงเป็น 26 จาก 28 หน้า — จำนวนหน้าทั้งหมดลดจาก 30 เป็น 28 เพราะยุบหน้าตั้งค่าผู้ดูแลระบบและหน้ารายงานเชิงวิเคราะห์เข้ากับหน้าที่ทำงานเดียวกัน
 \item ผู้ยืมยกเลิกคำขอของตนเองผ่าน backend ได้ ปุ่มจะซ่อนเองเมื่อเจ้าหน้าที่จัดของไปแล้ว ตามที่ backend บอก ไม่ได้คำนวณกฎซ้ำที่หน้าจอ
 \item ปุ่มออกจากระบบทุกอุปกรณ์ที่หน้าโปรไฟล์ ใช้ procedure ที่มีอยู่แล้วตั้งแต่ Sprint 1
 \item คำแปลภาษาไทยของข้อความผิดพลาดครอบคลุม 60 จาก 63 รหัสที่ backend ส่งได้ รอบก่อนครอบคลุม 29 จาก 59 ที่ขาดอยู่ 3 รหัสคือรหัสที่เพิ่งเพิ่มในรอบนี้
\end{itemize}

ข้อจำกัดที่เหลือมีสี่ข้อ ข้อแรกคือปุ่มส่งคำขอของผู้ยืมยังไม่เรียก backend และคำขอที่เพิ่งส่งไม่ปรากฏในรายการตามที่อธิบายในบทสรุป ข้อที่สองคือปุ่มขอต่ออายุของผู้ยืมยังบันทึกลง store ในเบราว์เซอร์ ทั้งที่ procedure ฝั่ง backend พร้อมแล้ว ข้อที่สามคือหน้าห้อง 3 หน้าและหน้าอุทธรณ์ 2 หน้ายังใช้ mock และหน้าส่งมอบกับหน้าส่งออกรายงานยังว่าง ข้อที่สี่คือยังใช้ไฟล์ contract ที่เขียนเอง ซึ่งประกาศ procedure ที่ไม่มีจริง 11 รายการ และไม่มี procedure ใหม่ของการต่ออายุเลย

\subsection{เครื่องมือและกระบวนการ}

งานตรวจเวอร์ชันของ zod ที่สองฝั่งใช้ร่วมกันเปลี่ยนเป็นงานบังคับใน CI แล้ว ปิดข้อจำกัดที่ค้างมาตั้งแต่รอบแรก และทั้งทีมใช้ Node เวอร์ชันเดียวกันที่ตรึงไว้ในไฟล์ของ repository

ข้อจำกัดที่เหลือเป็นข้อเดิมจากรอบก่อน คือฝั่ง frontend ยังไม่มี test runner และชุด AC check 44 ข้อยังต้องสั่งรันเอง ไม่ได้อยู่ใน CI

\subsection{เอกสารและการออกแบบ}

\begin{itemize}
 \item SDS ฉบับ 1.1 — นับทุกคอลัมน์สถานะใหม่จากโค้ดของวันรายงาน เพิ่มการออกแบบการต่ออายุและตาราง cron job และบันทึกข้อต่างระหว่างโค้ดกับ SRS ที่ต้องตัดสิน ใช้รูปแบบเดียวกับฉบับ 1.0 ที่ส่งไป
 \item หลักฐานการเรียก API จริงของวงจรยืม–คืน — 33 ครั้งบนฐานข้อมูลที่ตั้งขึ้นใหม่ ได้ผลตามคาด 29 ครั้ง อีก 4 ครั้งเป็นการทดสอบเชิงลบที่ตั้งใจให้ถูกปฏิเสธ และถูกปฏิเสธจริง
\end{itemize}

ข้อจำกัดที่เหลือคือยังไม่มีคู่มือผู้ใช้และคู่มือติดตั้ง ซึ่งอยู่ในแผน Sprint 5

\subsection{ระดับหลักฐานของสิ่งที่รายงานว่าเสร็จ}

\noindent\begin{tabularx}{\dimexpr\textwidth-\leftskip\relax}{@{} L{3.8cm} Y L{3.0cm} @{}}
\toprule
\hrow \thead{สิ่งที่ส่งมอบ} & \thead{หลักฐาน ณ วันรายงาน} & \thead{ระดับ} \\
\midrule
contract 86 procedure ใน 10 domain & ระบบ mount route ครบทุก procedure ตอนเริ่มทำงาน และ CI ผ่าน & พิสูจน์แล้ว \\
ยืนยันตัวตน session และ audit log & test ทำงานบนฐานข้อมูลจริงใน CI & พิสูจน์แล้ว \\
กฎเลือกเส้นทางอนุมัติและผู้ตัดสินการต่ออายุ & test 19 กรณีใน CI & พิสูจน์แล้ว \\
วงจรยืม–คืน การอนุมัติ และ notification & เรียก API ด้วยมือหนึ่งรอบเมื่อ 9 ก.ย. ครบทุกขั้น ไม่มี test & พิสูจน์แล้ว แต่ไม่อัตโนมัติ \\
กฎธุรกิจฝั่ง frontend & ชุด AC check 44 ข้อ ผ่านทั้งหมด ต้องสั่งรันเอง & พิสูจน์แล้ว แต่ไม่อัตโนมัติ \\
การต่ออายุฝั่ง service cron job และรายงาน & ไม่มี test และเข้า main หลังการเรียก API ด้วยมือ & เขียนเสร็จ ยังไม่พิสูจน์ \\
หน้าจอที่อ่านข้อมูลจริง 21 หน้า & typecheck ผ่าน ยังไม่มีการทดสอบหน้าจอ & เขียนเสร็จ ยังไม่พิสูจน์ \\
\bottomrule
\end{tabularx}

แถวที่ขยับขึ้นจากรอบก่อนคือวงจรยืม–คืน จาก ``เขียนเสร็จ ยังไม่พิสูจน์'' เป็น ``พิสูจน์แล้ว แต่ไม่อัตโนมัติ'' ส่วนงานที่เขียนใหม่ในรอบนี้เกือบทั้งหมดอยู่ในระดับต่ำสุด

% =============================================================================
\clearpage
\section{แผนการดำเนินงานถัดไป}
\label{sec:next}

\subsection{สี่วันแรกของ Sprint 4 (11 – 14 กันยายน 2569) — งานค้างของ Sprint 3}

งานในช่วงนี้ทั้งหมดเป็นการต่อสิ่งที่มีอยู่แล้วเข้าหากัน ไม่มีความสามารถใหม่ ลำดับเรียงตามสิ่งที่ขวางผู้ใช้ก่อน

\noindent\begin{tabularx}{\dimexpr\textwidth-\leftskip\relax}{@{} L{1.5cm} Y L{2.1cm} @{}}
\toprule
\hrow \thead{กำหนด} & \thead{งาน} & \thead{ผู้รับผิดชอบ} \\
\midrule
11 ก.ย. & ต่อปุ่มส่งคำขอของผู้ยืมเข้ากับ procedure สร้างคำขอ แก้ข้อถดถอยที่คำขอหายจากรายการไปพร้อมกัน & FE-dev \\
11 ก.ย. & เพิ่มคำแปลไทยของ 3 รหัสข้อผิดพลาดใหม่ & FE-dev \\
12 ก.ย. & หน้าส่งมอบของเจ้าหน้าที่ ใช้ procedure จัดของ สลับชิ้น และยืนยันการรับที่มีอยู่แล้ว & FE-lead \\
12 ก.ย. & ต่อปุ่มขอต่ออายุของผู้ยืม และคิวพิจารณาการต่ออายุของเจ้าหน้าที่และอาจารย์ เข้ากับ procedure ที่มีอยู่แล้ว & FE-lead \\
12 ก.ย. & ตัดสินข้อต่างเรื่องการต่ออายุระหว่างโค้ดกับ FR-RNW-04 และ FR-RNW-06 แล้วแก้ฝั่งใดฝั่งหนึ่ง & SA \\
13 ก.ย. & test ที่เดินวงจรยืม–คืนตั้งแต่สร้างคำขอจนรับคืนบนฐานข้อมูลจริง ให้ทำงานใน CI & BE-lead \\
13 ก.ย. & ติดตั้ง test runner ฝั่ง frontend และย้ายชุด AC check 44 ข้อเข้า CI & QA-lead \\
14 ก.ย. & เดินวงจรยืม–คืนครบเส้นทางปกติผ่านหน้าจอจริง ตั้งแต่ผู้ยืมสร้างคำขอจนเจ้าหน้าที่รับคืน & PM \\
14 ก.ย. & เก็บไฟล์ contract ที่สร้างจากโค้ดเข้า repository และรับไปใช้แทนไฟล์ที่เขียนเอง & BE-lead \newline FE-lead \\
\bottomrule
\end{tabularx}

\subsection{ส่วนที่เหลือของ Sprint 4 และ Sprint 5}

\noindent\begin{tabularx}{\dimexpr\textwidth-\leftskip\relax}{@{} L{2.9cm} Y @{}}
\toprule
\hrow \thead{Sprint} & \thead{เป้าหมายและเกณฑ์ ``เสร็จ''} \\
\midrule
\textbf{Sprint 4} \newline 15 – 24 ก.ย. & การจองสถานที่ฝั่ง backend และต่อหน้าห้อง 3 หน้า · การอุทธรณ์ฝั่ง backend และต่อหน้าอุทธรณ์ 2 หน้า · ตะกร้าพร้อมการแตกอนุมัติตามหน่วยงานเจ้าของ · การจองล่วงหน้า · cron job ที่เหลือ 3 งาน · หน้าส่งออกรายงาน · การค้นหาและการรองรับมือถือ \newline เกณฑ์เสร็จ — ไม่เหลือหน้าว่างและไม่เหลือ mock ในระบบ \newline หมายเหตุ — การต่ออายุฝั่ง backend เสร็จแล้ว เป็น Sprint สุดท้ายที่เพิ่มความสามารถใหม่ได้ \\
\textbf{Sprint 5} \newline 25 ก.ย. – 1 ต.ค. & หยุดเพิ่มความสามารถและเตรียมนำเสนอ — ทดสอบถดถอยเต็มรูปแบบ · ทดสอบสมรรถนะ · นำระบบขึ้นใช้งานจริงพร้อมย้ายที่เก็บไฟล์รูปไปไว้ภายนอกและตรวจสิทธิ์การเปิดรูป · ปรับ SRS และ SDS ให้ตรงกับระบบที่ส่งมอบ · คู่มือผู้ใช้และคู่มือติดตั้ง · ซ้อมนำเสนอ \\
\bottomrule
\end{tabularx}

\subsection{การประเมินกำหนดเสร็จ}

คณะทำงานประเมินว่าโครงการยังเสร็จได้ภายใน 1 ตุลาคม 2569 ด้วยขอบเขตครบตามข้อเสนอโครงการ แต่ความมั่นใจลดลงจากรอบก่อน เหตุผลที่ยังคงกำหนดเดิมคือความสามารถของ backend เดินหน้ากว่าแผน การต่ออายุซึ่งเป็นงานของ Sprint 4 เสร็จไปแล้ว และงานค้างทั้งหมดเป็นงานต่อเชื่อมที่ประเมินเวลาได้ เหตุผลที่ความมั่นใจลดลงคืองานต่อเชื่อมชุดเดียวกันถูกเลื่อนมาสองรอบติดกัน ปัญหาจึงไม่ใช่ความยากของงาน แต่เป็นการที่งานประเภทนี้ถูกแซงด้วยงานสร้างของใหม่ซ้ำ ๆ

การประเมินข้างต้นตั้งอยู่บนจุดวัดสี่ข้อถัดไป ถ้าข้อ 1 หรือข้อ 2 ไม่ผ่าน คณะทำงานจะลดขอบเขตของ Sprint 4 ในส่วนตะกร้าข้ามหน่วยงานและการจองล่วงหน้าก่อน เพื่อรักษาเวลาของ Sprint 5 ไว้สำหรับการทดสอบ

\begin{enumerate}
 \item ผู้ยืมส่งคำขอผ่านหน้าจอแล้วคำขอถูกเขียนลงฐานข้อมูล และหน้าส่งมอบของเจ้าหน้าที่ใช้งานได้ ภายใน 12 กันยายน 2569
 \item เดินวงจรยืม–คืนครบเส้นทางปกติผ่านหน้าจอจริง และมี test ของวงจรนี้ทำงานใน CI ภายใน 14 กันยายน 2569
 \item ฝั่ง frontend ใช้ไฟล์ contract ที่สร้างจากโค้ด ภายใน 14 กันยายน 2569
 \item ไม่เหลือ mock และหน้าว่างในระบบ ภายในวันสิ้นสุด Sprint 4 คือ 24 กันยายน 2569
\end{enumerate}

% =============================================================================
\clearpage
\section*{ภาคผนวก — ค่าที่ใช้พล็อตกราฟ}

ค่าเป้าหมายรายมิติ ณ วันสิ้นสุดของแต่ละ Sprint คูณด้วยน้ำหนักของมิตินั้นแล้วรวมกัน คอลัมน์ ``รวม'' คือค่าที่พล็อตบนเส้นแผนในรูปที่ \ref{fig:scurve} แถวที่มีเครื่องหมายดอกจันคือวันรายงานรอบนี้ ซึ่งตรงกับขอบ Sprint 3 ตารางนี้เป็นตารางเดียวกับรายงานครั้งที่ 2 ไม่มีการปรับแผน

\noindent\begin{tabularx}{\dimexpr\textwidth-\leftskip\relax}{@{} L{3.2cm} L{1.0cm} L{1.0cm} L{1.0cm} L{1.0cm} L{1.0cm} L{1.0cm} Y @{}}
\toprule
\hrow \thead{ขอบ Sprint} & \thead{BE Des} & \thead{FE Des} & \thead{BE Dev} & \thead{FE Dev} & \thead{Test} & \thead{Doc} & \thead{รวม} \\
\hrow \thead{} & \thead{10\%} & \thead{10\%} & \thead{30\%} & \thead{25\%} & \thead{15\%} & \thead{10\%} & \thead{100\%} \\
\midrule
เริ่ม 24 ก.ค. & 0 & 0 & 0 & 0 & 0 & 10 & 1 \\
S0 · 30 ก.ค. & 45 & 30 & 5 & 5 & 0 & 20 & 12 \\
S1 · 13 ส.ค. & 88 & 75 & 25 & 42 & 10 & 65 & 42 \\
S2 · 27 ส.ค. & 95 & 85 & 60 & 65 & 35 & 70 & 65 \\
$*$ S3 · 10 ก.ย. & 100 & 92 & 78 & 78 & 52 & 76 & 78 \\
S4 · 24 ก.ย. & 100 & 100 & 97 & 97 & 80 & 88 & 94 \\
S5 · 1 ต.ค. & 100 & 100 & 100 & 100 & 100 & 100 & 100 \\
\bottomrule
\end{tabularx}

ค่าผลจริง ณ วันรายงาน คือ 97, 95, 77, 74, 35 และ 88 ตามลำดับมิติเดียวกัน ถ่วงน้ำหนักแล้วได้ 74.9 ปัดเป็น \actualpct

ค่าผลจริงรอบก่อนหน้าที่อยู่บนเส้นผลจริง คือ 52\% ณ 27 สิงหาคม และ 70\% ณ 4 กันยายน ตามที่ส่งในรายงานครั้งที่ 1 และครั้งที่ 2 ไม่มีการแก้ย้อนหลัง

\end{document}
